\documentclass[11pt]{article}
\usepackage[margin=25mm]{geometry}
\usepackage{fontspec}
\setmainfont{Latin Modern Roman}
\usepackage{amsmath,amssymb,mathtools}
\usepackage{unicode-math}
\setmathfont{Latin Modern Math}
\usepackage{xurl,hyperref}
\hypersetup{colorlinks=true,urlcolor=blue}
\setcounter{secnumdepth}{0}
\setlength{\parindent}{0pt}
\setlength{\parskip}{5pt}
\setlength{\emergencystretch}{3em}
\newcommand{\sref}[2]{\hyperlink{source:#1:#2}{[S#2]}}
\newcommand{\cataloguescope}[1]{\paragraph{Recorded target.}#1}
\newcommand{\R}{\mathbb R}
\newcommand{\C}{\mathbb C}
\newcommand{\N}{\mathbb N}
\newcommand{\Z}{\mathbb Z}
\newcommand{\Q}{\mathbb Q}
\newcommand{\F}{\mathbb F}
\newcommand{\eps}{\varepsilon}
\begin{document}
% UnsolvedMath ID 8100006; source catalogue code AMR-080-0006
\setcounter{section}{8100005}
\section{numerical solution of integrable systems}\label{8100006}
{\small UnsolvedMath ID 8100006 · Mathematical Physics}
\subsection{Definitions and mathematical statement}

\paragraph{Shared notation.}$\mathbb N=\{1,2,\ldots\}$; the standard number fields have their usual meanings. All additional conventions are stated locally.


\paragraph{Definitions and notation.}The inverse-scattering forward map sends initial data to a reflection coefficient and, where present, discrete eigenvalues and norming constants. The inverse map is a matrix Riemann--Hilbert problem. Small dispersion introduces $\varepsilon\downarrow0$ in the forward spectral problem. The Airy operator on $L^2(s,\infty)$ has kernel $K(\lambda,\mu)=[\operatorname{Ai}(\lambda)\operatorname{Ai}'(\mu)-\operatorname{Ai}'(\lambda)\operatorname{Ai}(\mu)]/(\lambda-\mu)$, with diagonal defined by continuity.

\paragraph{Catalogue wording retained for traceability.}
Apply and extend the Olver--Deconinck--Trogdon numerical Riemann--Hilbert framework to integrable systems, addressing both forward scattering and inverse reconstruction. In particular, solve the non-self-adjoint forward problem for focusing NLS with general smooth initial data and reconstruct the solution in the small-dispersion limit $\varepsilon=\hbar\downarrow0$. Analyze Airy-operator eigenfunctions as $s\to\infty$ and compute them reliably for finite $s$. Continue the Painlevé resource project collecting algebraic, analytic and numerical information, with reliable computation of representative solutions.
\paragraph{Scope and source qualification.}The original is explicitly an open programme. The catalogue’s cut-off ending is restored from primary Problem 6. Known sine-kernel/prolate-spheroidal work is context rather than an extra claim of unresolved status.
\subsection{Short English statement}
Develop reliable numerical inverse scattering, resolve general focusing small-dispersion data, and analyze Airy-operator eigenfunctions.
\subsection{Research attempt}
\paragraph{Model, date and investigated scope.}
OpenAI GPT-6 Astra (configured model \texttt{gpt-6-astra}), ultra reasoning; generated 11 October 2026. Derived Airy-kernel positivity, diagonal formula and a positivity-preserving finite-factor approximation principle. Eigenfunction asymptotics, certified computation and focusing inverse-scattering tasks remain unresolved.

\paragraph{Sources and literature check.}
Read the supplied complete definition, catalogue wording, scope qualifications and source list for this canonical ID. The original-source citations are inherited as supplied references; no fresh primary-source or novelty search was performed for this attempt. No current open-status claim is inferred from this calculation.

\paragraph{A positive representation of the Airy kernel.}Let \(a=\operatorname{Ai}\), satisfying \(a''(x)=xa(x)\) and decaying at positive infinity. Define \(W(t)=a(x+t)a'(y+t)-a'(x+t)a(y+t)\). Then \(W'(t)=(y-x)a(x+t)a(y+t)\). Integration from zero to infinity, using the decay of \(a,a'\), gives
\[K(x,y)=\frac{a(x)a'(y)-a'(x)a(y)}{x-y}=\int_0^\infty a(x+t)a(y+t)\,dt.
\]
Consequently for compactly supported \(f\),
\[\langle f,Kf\rangle=\int_0^\infty\left|\int_s^\infty a(x+t)f(x)dx\right|^2dt\ge0.
\]
The diagonal is \(K(x,x)=a'(x)^2-xa(x)^2\), obtained by differentiating numerator and denominator with respect to \(y\), including the minus sign in the denominator derivative.
\paragraph{What this offers for computation.}If an integral quadrature with positive weights approximates the displayed factorization by a finite matrix \(B\), the approximation \(BB^*\) is positive semidefinite automatically. This preserves positivity but supplies no error bound. Small operator norm also does not by itself identify a normalized eigenfunction when adjacent eigenvalues are nearly equal; the needed perturbation control must be small relative to their gap.
\paragraph{First gap.}No asymptotic eigenfunction profile as \(s\to\infty\), certified quadrature error or finite-\(s\) implementation is produced. Establishing a uniform rescaled kernel expansion with controlled spectral gaps is missing even for the Airy subproblem. The non-self-adjoint general-data focusing forward problem, inverse reconstruction at small dispersion and broader Painleve software resource remain unaddressed. This exact kernel identity is a standard analytic benchmark, not a claimed completion of the numerical programme.

\paragraph{Verification and adversarial audit.}
Checked Wronskian derivative, integral sign, diagonal limit, positive quadratic form and distinction between positivity and error certification. This was an LLM adversarial audit, not independent human verification; no theorem prover was used. Prompts A then B applied to this ID: Checked Wronskian derivative, integral sign, diagonal limit, positive quadratic form and distinction between positivity and error certification. Scope remains limited: Obtain uniform rescaled kernel and spectral-gap estimates sufficient for certified Airy eigenfunction computation, separately from the nonnormal scattering problem.

Coordinator checked the Wronskian derivative, kernel integral identity, positivity and diagonal sign. This additional check was model-assisted, not independent human verification.

\paragraph{Tools, process and contributions.}
Python 3 and local shell used to read supplied files and serialize this manuscript; mathematics derived symbolically in the text, with no numerical experiment or formal proof assistant. Contribution prompt A was followed by prompt B. GPT-6 Astra with ultra reasoning drafted the research and performed its stated self-audit; any additional coordinator contribution is identified in the verification and audit. A GPT-6-based Codex coordinator assembled the manuscript and checked repository integration; its exact variant and reasoning configuration were not exposed. No human mathematical contribution, independent human proof check or formal proof certification is supplied. Complete derivations and the consulted source references are included here; private reasoning transcripts are not published.

\paragraph{Final status and first remaining gap.}
\textbf{UNRESOLVED}. Obtain uniform rescaled kernel and spectral-gap estimates sufficient for certified Airy eigenfunction computation, separately from the nonnormal scattering problem.

\section{Completion Estimate}
\noindent\textbf{COMPLETION: 0\%}
The calculation is an elementary benchmark or conditional reduction; it establishes no claimed new portion of the complete source target.

\subsection{Sources}
\begin{itemize}
\item[\hypertarget{source:8100006:1}{S1}] ULAM.AI and the UnsolvedMath Contributors, UnsolvedMath, record 8100006 (AMR-080-0006). Catalogue curation under CC BY 4.0; definitions and English exposition added locally.
\url{https://huggingface.co/datasets/ulamai/UnsolvedMath}.
\item[\hypertarget{source:8100006:2}{S2}] Percy Deift, Some Open Problems in Random Matrix Theory and the Theory of Integrable Systems. II, SIGMA 13 (2017), 016, Problem 6.
\url{https://arxiv.org/abs/1703.04931}.
\item[\hypertarget{source:8100006:3}{S3}] Primary full-text formulation and 2017 historical progress discussion.
\url{https://arxiv.org/html/1703.04931v1}.
\end{itemize}
\label{end:8100006}
\end{document}
